% Lecture 10 - Page 3 (6-XII-1994).
% Independent fragment for % Lecture 10 - Page 3 (6-XII-1994).
% Independent fragment for % Lecture 10 - Page 3 (6-XII-1994).
% Independent fragment for % Lecture 10 - Page 3 (6-XII-1994).
% Independent fragment for \input{courses/pages/curs10-pagina3-en}.
% Continuation of the proof begun on page 2.

\subsection*{Lecture 10 - Page 3}

\noindent (2) \quad
$N\trianglelefteq G\ \Longleftrightarrow\ axa^{-1}\in N$,
\quad $(\forall)\,a\in G$, $(\forall)\,x\in N$.

\medskip
\noindent By the definition of $N_G(H)$,
\[
axa^{-1}\in H, \qquad
(\forall)\,a\in N_G(H),\quad (\forall)\,x\in H.
\]
\noindent $\Longrightarrow\ H\trianglelefteq N_G(H)$.
We may consider the group $N_G(H)/H$.

\bigskip
% An isolated/cancelled mark precedes the dividing line in the manuscript.
\hrule
\bigskip

\noindent Consider an arbitrary set $A$, $|A|<\infty$.

\medskip
\noindent Let $z_1,\ldots,z_q\in A$ be a transversal of the equivalence
relation associated with the action of $G$ on $A$, ``$\sim_G$''.

\[
\left\{\begin{aligned}
&i\neq j\ \Longrightarrow\ z_i\not\sim_G z_j,\\
&(\forall)\,x\in A,\ (\exists)\,i\ \text{such that}\ x\sim_G z_i.
\end{aligned}\right.
\]

\noindent Under these conditions we have:
\[
A=\operatorname{Orb}(z_1)\cup\operatorname{Orb}(z_2)
\cup\cdots\cup\operatorname{Orb}(z_q)
\qquad (\text{disjoint union}).
\]

\[
\Longrightarrow\quad
|A|=\sum_{i=1}^{q}[G:\operatorname{Stab}_G(z_i)].
\]
\noindent \underline{Class equation for the action of $G$ on $A$}.

\bigskip
\noindent In particular, if $A=G$, considering the action by conjugation
of $G$ on itself ($G$ a finite group),
\[
\Longrightarrow\quad
\operatorname{ord}(G)=\sum_{i=1}^{q}[G:C_G(z_i)]
\]
\noindent - the class equation for conjugacy classes of elements.

\bigskip
\noindent Let $x\in G$ be such that
\[
C_G(x)=G\quad\Longleftrightarrow\quad
|\operatorname{orb}_G(x)|=1,
\]
\noindent i.e. $ax=xa$, $(\forall)\,a\in G$.

\bigskip
\noindent We denote:
\[
Z_G:=\{x\in G\mid ax=xa,\ (\forall)\,a\in G\}.
\]

\noindent $Z_G\trianglelefteq G$.
\quad $Z_G$ is called \underline{the centre of the group $G$}.

\medskip
\noindent $x,y\in Z_G\ \stackrel{?}{\Longrightarrow}\ xy^{-1}\in Z_G$;
\quad and if $x\in Z_G$, $a\in G$,
\[
\stackrel{?}{\Longrightarrow}\ axa^{-1}\in Z_G
\qquad (axa^{-1}=xaa^{-1}=x).
\]

\bigskip
\noindent $Z_G$ is contained in every transversal; in particular,
$Z_G\subseteq\{z_1,\ldots,z_q\}=T$.

\[
\operatorname{ord}(G)=\operatorname{ord}(Z_G)
+\sum_{\substack{z_i\in T\\z_i\notin Z_G}}[G:C_G(z_i)].
\]
% For central z_i, [G:C_G(z_i)]=1 (handwritten annotation "1").
.
% Continuation of the proof begun on page 2.

\subsection*{Lecture 10 - Page 3}

\noindent (2) \quad
$N\trianglelefteq G\ \Longleftrightarrow\ axa^{-1}\in N$,
\quad $(\forall)\,a\in G$, $(\forall)\,x\in N$.

\medskip
\noindent By the definition of $N_G(H)$,
\[
axa^{-1}\in H, \qquad
(\forall)\,a\in N_G(H),\quad (\forall)\,x\in H.
\]
\noindent $\Longrightarrow\ H\trianglelefteq N_G(H)$.
We may consider the group $N_G(H)/H$.

\bigskip
% An isolated/cancelled mark precedes the dividing line in the manuscript.
\hrule
\bigskip

\noindent Consider an arbitrary set $A$, $|A|<\infty$.

\medskip
\noindent Let $z_1,\ldots,z_q\in A$ be a transversal of the equivalence
relation associated with the action of $G$ on $A$, ``$\sim_G$''.

\[
\left\{\begin{aligned}
&i\neq j\ \Longrightarrow\ z_i\not\sim_G z_j,\\
&(\forall)\,x\in A,\ (\exists)\,i\ \text{such that}\ x\sim_G z_i.
\end{aligned}\right.
\]

\noindent Under these conditions we have:
\[
A=\operatorname{Orb}(z_1)\cup\operatorname{Orb}(z_2)
\cup\cdots\cup\operatorname{Orb}(z_q)
\qquad (\text{disjoint union}).
\]

\[
\Longrightarrow\quad
|A|=\sum_{i=1}^{q}[G:\operatorname{Stab}_G(z_i)].
\]
\noindent \underline{Class equation for the action of $G$ on $A$}.

\bigskip
\noindent In particular, if $A=G$, considering the action by conjugation
of $G$ on itself ($G$ a finite group),
\[
\Longrightarrow\quad
\operatorname{ord}(G)=\sum_{i=1}^{q}[G:C_G(z_i)]
\]
\noindent - the class equation for conjugacy classes of elements.

\bigskip
\noindent Let $x\in G$ be such that
\[
C_G(x)=G\quad\Longleftrightarrow\quad
|\operatorname{orb}_G(x)|=1,
\]
\noindent i.e. $ax=xa$, $(\forall)\,a\in G$.

\bigskip
\noindent We denote:
\[
Z_G:=\{x\in G\mid ax=xa,\ (\forall)\,a\in G\}.
\]

\noindent $Z_G\trianglelefteq G$.
\quad $Z_G$ is called \underline{the centre of the group $G$}.

\medskip
\noindent $x,y\in Z_G\ \stackrel{?}{\Longrightarrow}\ xy^{-1}\in Z_G$;
\quad and if $x\in Z_G$, $a\in G$,
\[
\stackrel{?}{\Longrightarrow}\ axa^{-1}\in Z_G
\qquad (axa^{-1}=xaa^{-1}=x).
\]

\bigskip
\noindent $Z_G$ is contained in every transversal; in particular,
$Z_G\subseteq\{z_1,\ldots,z_q\}=T$.

\[
\operatorname{ord}(G)=\operatorname{ord}(Z_G)
+\sum_{\substack{z_i\in T\\z_i\notin Z_G}}[G:C_G(z_i)].
\]
% For central z_i, [G:C_G(z_i)]=1 (handwritten annotation "1").
.
% Continuation of the proof begun on page 2.

\subsection*{Lecture 10 - Page 3}

\noindent (2) \quad
$N\trianglelefteq G\ \Longleftrightarrow\ axa^{-1}\in N$,
\quad $(\forall)\,a\in G$, $(\forall)\,x\in N$.

\medskip
\noindent By the definition of $N_G(H)$,
\[
axa^{-1}\in H, \qquad
(\forall)\,a\in N_G(H),\quad (\forall)\,x\in H.
\]
\noindent $\Longrightarrow\ H\trianglelefteq N_G(H)$.
We may consider the group $N_G(H)/H$.

\bigskip
% An isolated/cancelled mark precedes the dividing line in the manuscript.
\hrule
\bigskip

\noindent Consider an arbitrary set $A$, $|A|<\infty$.

\medskip
\noindent Let $z_1,\ldots,z_q\in A$ be a transversal of the equivalence
relation associated with the action of $G$ on $A$, ``$\sim_G$''.

\[
\left\{\begin{aligned}
&i\neq j\ \Longrightarrow\ z_i\not\sim_G z_j,\\
&(\forall)\,x\in A,\ (\exists)\,i\ \text{such that}\ x\sim_G z_i.
\end{aligned}\right.
\]

\noindent Under these conditions we have:
\[
A=\operatorname{Orb}(z_1)\cup\operatorname{Orb}(z_2)
\cup\cdots\cup\operatorname{Orb}(z_q)
\qquad (\text{disjoint union}).
\]

\[
\Longrightarrow\quad
|A|=\sum_{i=1}^{q}[G:\operatorname{Stab}_G(z_i)].
\]
\noindent \underline{Class equation for the action of $G$ on $A$}.

\bigskip
\noindent In particular, if $A=G$, considering the action by conjugation
of $G$ on itself ($G$ a finite group),
\[
\Longrightarrow\quad
\operatorname{ord}(G)=\sum_{i=1}^{q}[G:C_G(z_i)]
\]
\noindent - the class equation for conjugacy classes of elements.

\bigskip
\noindent Let $x\in G$ be such that
\[
C_G(x)=G\quad\Longleftrightarrow\quad
|\operatorname{orb}_G(x)|=1,
\]
\noindent i.e. $ax=xa$, $(\forall)\,a\in G$.

\bigskip
\noindent We denote:
\[
Z_G:=\{x\in G\mid ax=xa,\ (\forall)\,a\in G\}.
\]

\noindent $Z_G\trianglelefteq G$.
\quad $Z_G$ is called \underline{the centre of the group $G$}.

\medskip
\noindent $x,y\in Z_G\ \stackrel{?}{\Longrightarrow}\ xy^{-1}\in Z_G$;
\quad and if $x\in Z_G$, $a\in G$,
\[
\stackrel{?}{\Longrightarrow}\ axa^{-1}\in Z_G
\qquad (axa^{-1}=xaa^{-1}=x).
\]

\bigskip
\noindent $Z_G$ is contained in every transversal; in particular,
$Z_G\subseteq\{z_1,\ldots,z_q\}=T$.

\[
\operatorname{ord}(G)=\operatorname{ord}(Z_G)
+\sum_{\substack{z_i\in T\\z_i\notin Z_G}}[G:C_G(z_i)].
\]
% For central z_i, [G:C_G(z_i)]=1 (handwritten annotation "1").
.
% Continuation of the proof begun on page 2.

\subsection*{Lecture 10 - Page 3}

\noindent (2) \quad
$N\trianglelefteq G\ \Longleftrightarrow\ axa^{-1}\in N$,
\quad $(\forall)\,a\in G$, $(\forall)\,x\in N$.

\medskip
\noindent By the definition of $N_G(H)$,
\[
axa^{-1}\in H, \qquad
(\forall)\,a\in N_G(H),\quad (\forall)\,x\in H.
\]
\noindent $\Longrightarrow\ H\trianglelefteq N_G(H)$.
We may consider the group $N_G(H)/H$.

\bigskip
% An isolated/cancelled mark precedes the dividing line in the manuscript.
\hrule
\bigskip

\noindent Consider an arbitrary set $A$, $|A|<\infty$.

\medskip
\noindent Let $z_1,\ldots,z_q\in A$ be a transversal of the equivalence
relation associated with the action of $G$ on $A$, ``$\sim_G$''.

\[
\left\{\begin{aligned}
&i\neq j\ \Longrightarrow\ z_i\not\sim_G z_j,\\
&(\forall)\,x\in A,\ (\exists)\,i\ \text{such that}\ x\sim_G z_i.
\end{aligned}\right.
\]

\noindent Under these conditions we have:
\[
A=\operatorname{Orb}(z_1)\cup\operatorname{Orb}(z_2)
\cup\cdots\cup\operatorname{Orb}(z_q)
\qquad (\text{disjoint union}).
\]

\[
\Longrightarrow\quad
|A|=\sum_{i=1}^{q}[G:\operatorname{Stab}_G(z_i)].
\]
\noindent \underline{Class equation for the action of $G$ on $A$}.

\bigskip
\noindent In particular, if $A=G$, considering the action by conjugation
of $G$ on itself ($G$ a finite group),
\[
\Longrightarrow\quad
\operatorname{ord}(G)=\sum_{i=1}^{q}[G:C_G(z_i)]
\]
\noindent - the class equation for conjugacy classes of elements.

\bigskip
\noindent Let $x\in G$ be such that
\[
C_G(x)=G\quad\Longleftrightarrow\quad
|\operatorname{orb}_G(x)|=1,
\]
\noindent i.e. $ax=xa$, $(\forall)\,a\in G$.

\bigskip
\noindent We denote:
\[
Z_G:=\{x\in G\mid ax=xa,\ (\forall)\,a\in G\}.
\]

\noindent $Z_G\trianglelefteq G$.
\quad $Z_G$ is called \underline{the centre of the group $G$}.

\medskip
\noindent $x,y\in Z_G\ \stackrel{?}{\Longrightarrow}\ xy^{-1}\in Z_G$;
\quad and if $x\in Z_G$, $a\in G$,
\[
\stackrel{?}{\Longrightarrow}\ axa^{-1}\in Z_G
\qquad (axa^{-1}=xaa^{-1}=x).
\]

\bigskip
\noindent $Z_G$ is contained in every transversal; in particular,
$Z_G\subseteq\{z_1,\ldots,z_q\}=T$.

\[
\operatorname{ord}(G)=\operatorname{ord}(Z_G)
+\sum_{\substack{z_i\in T\\z_i\notin Z_G}}[G:C_G(z_i)].
\]
% For central z_i, [G:C_G(z_i)]=1 (handwritten annotation "1").
